\homework{6}{Mon 13 Feb 2023 17:07}{}

\begin{itemize}
	\item Suppose $C \subset M$ has the property that any cover of $C$ by balls has a finite subcover. Prove that $C$ is compact.
\begin{proof}
	Let $\mathcal{C}$ be a open cover of $C$. For each  $A \in \mathcal{C}$, define the collection of open sets $\mathcal{C}_A$ as
	\[
		\mathcal{C}_A := \{B_{r(a)}(a): a \in A\} 
	.\] 
	where $r(a)$ is some positive real number such that $B_{r(a)}(a) \subset A$. Such $r$ always exists by openness of $A$. Since each element in $\mathcal{C}_A$ is contained in $A$, we have $\cup \mathcal{C}_A \subset A$. Also each $a \in A$ is contained in $B_{r(a)}(a)$, so $A \subset \cup \mathcal{C}_A$. Finally
	\[
		\mathcal{D}:=\cup_{A \in \mathcal{C}} \mathcal{C}_{A}	.\] 
		covers $\cup \mathcal{C}$ by open balls, thus also covers $C$. Using the assumption we have a finite subset $\mathcal{D}' \subset \mathcal{D}$ that covers $C$. Let $\tau: \mathcal{D} \to \mathcal{C}$ which assigns to each $B_{r(a)}(a) \in \mathcal{D}$ some $A \in \mathcal{C}$ such that $a \in A$. This induces a corresponding subset of $\mathcal{C}$ :
		\[
		\mathcal{C}' = \tau(\mathcal{D}') 
		.\] 
		Since $D'$ is finite, $ \mathcal{C}'$ is finite. Also, for each $x \in C$ we have $x \in B_{r(a)}(a)$ for some $B_{r(a)}(a) \in \mathcal{D}'$. Furthermore, $ B_{r(a)}(a) \subset \tau(B_{r(a)}(a))  \in \mathcal{C}'$ by definition of $r$ and $\tau$. Hence for each $x \in C$ there exists some set in $\mathcal{C}'$ that contains $x$. Therefore $\mathcal{C}'$ is the finite subcover we want. Since $\mathcal{C}$ is arbitrary open cover, $C$ is compact.
\end{proof}

\item A map $\varphi: M \to  M$ is called an isometry if $\rho(\varphi(x), \varphi(y)) = \rho(x,y)$ for all $x, y \in M$. Prove that $I = \{\varphi: M \to  M \text{ is an isometry}\} \subset \text{Map}(M, M)$ is compact.
\begin{proof}
It suffices to show $I$ is totally bounded and complete. We first show $I$ is totally bounded.

Fix arbitrary $\varepsilon>0$. Using compactness of $M$, choose a finite collection of points $A \subset M$ such that the balls $B_{\varepsilon}(a_i)$ cover $M$ where $a_i \in A, i = 1\ldots n$. In other words, let $A$ be a $\varepsilon$-net of $M$. Now define the set $X \subset A^A$ be
\[
	X = \{f \in A^A: \exists g_f \in I, \forall a \in A, \left| g_f(a) - f(a) \right| < \varepsilon\} 
.\] 
$X$ is finite since  $A^A$ has finite cardinality $n^n$. Now for each $f \in X$, choose $g_f \in I$ as in the definition of $X$. We claim that for every isometry $h \in I$ there exists $f \in X$ such that $\tau(g_f, h) \le  4 \varepsilon$.

Fix $h \in I$. For each $a_i \in A$, let $f(a_i) \in A$ be such that $\rho\left( h(a_i) , f(a_i) \right) < \varepsilon$. We know $f \in X$ exists by total-boundedness of $M$. Fix arbitrary $x \in M$, then there exists  $a_i$ such that $\rho(x, a_i) < \varepsilon$. 

By property of isometry, $\rho\left( h(x), h(a_i) \right) <\varepsilon$ and $\rho\left( g_f(x), g_f(a_i) \right) <\varepsilon$. By definition of $g_f$, $\rho\left(g_f(a_i), f(a_i)\right)<\varepsilon$.
Now by triangle inequality
\begin{align*}
	\rho\left(g_f(x), h(x)\right) &\le \rho\left( g_f(x), g_f(a_i) \right) + \rho\left( g_f(a_i) , f(a_i) \right) \\
				      &\qquad +\rho\left( f(a_i), h(a_i) \right) +\rho\left(h(a_i), h(x)\right) \\
				      &< 4 \varepsilon
.\end{align*} 

This is true for all $x$, so $\tau(g_f, h) \le 4 \varepsilon$. This implies total boundedness since the collections of balls centered at $\{g_f: f \in X\} $ with radius $4 \varepsilon$ would cover $I$. 

Finally we show completeness. Consider some Cauchy sequence $f_n \in I$.  Since $M$ compact, $M$ is complete, so $\text{Map}(M, M)$ is complete. Hence $f_n$ uniformly converges to some $f \in \text{Map}(M, M)$. It remains to show $f$ is an isometry. 
\[
	\rho(f(x), f(y)) = \rho(\lim f_n(x), \lim f_n(y)) = \lim \rho(f_n(x), f_n(y)) = \lim \rho(x,y) = \rho(x,y)
.\] 

We could exchange the limit with $\rho$ because $\rho$ is continuous.
\end{proof}
\item Suppose $f: \mathbb{R}\to \mathbb{R}$ has only finitely many points of discontinuity. Prove that there is a sequence $f_k: \mathbb{R}\to \mathbb{R}$ of continuous functions with $f_k \to f$ as $k \to \infty $.
\begin{proof}
	Let $x_1<  \ldots < x_n $ be all the points of discontinuity and let $\varepsilon = \max_i x_i - x_{i-1} > 0$. For $k > \frac{2}{\varepsilon}$ we define $f_k$ to be
	\[
		f_k(x) = \begin{cases}
			f(x) & |x-x_i| \ge  \frac{1}{k} \text{ for all $i$}\\
			f(x) & x = x_i \text{ for some $i$} \\
			\text{ linear in between }
		\end{cases}
	.\] 
	by choosing large enough $k$, $f_ k$ is well-defined. We verify $f_k$ is continuous. Fix arbitrary $x \in \mathbb{R}$. If $x = x_i$ for some $i$ then $f$ is continuous at $x$ since it is linear interpolation at $x$. If not, then for large enough $k$ we have $x$ is at least $\frac{1}{k}$-far from any $x_i$, so $f_k(x)$ eventually coincides with $f(x)$ which is continuous at $x$. Moreover, by the argument above, we also see that $f_k(x)$ eventually coincides with $f(x)$ for all $x \in \mathbb{R}$. Hence $f_k$ pointwise converges of $f$.
\end{proof}
\item Prove that a function $\varphi: M \to P$ is continuous at $a \in M$ iff for any sequence $x_n \in M$ with  $\lim_{n \to \infty} x_n = a$ we have $\varphi(a) = \lim_{n \to \infty} \varphi(x_n)$.
\begin{proof}
	Suppose that  $\varphi$ is continuous at $a$. Take a sequence $x_n \in M$ with limit $a$. Now by continuity of the function, for all $\varepsilon>0$ there exists $\delta>0$ such that $\left| f(x) - f(a) \right| <\varepsilon$ whenever $\left| x - a \right| <\delta$. By convergence of $x_n$, $\left| x_n - a \right| <\delta$ is satisfied by sufficiently large $n \ge N$. Thus $\left| f(x_n) - f(a) \right| <\varepsilon$ for sufficiently large $n$. This proves $f(x_n) \to f(a)$.

	Now suppose $f$ preserves limits of all sequences in $M$, we want to prove $f$ continuous. Suppose $f$ discontinuous at $a$, then there exists  $\varepsilon>0$ such that for any $\delta>0$ there exists $x \in M$ such that $\left| x - a \right| < \delta$ and yet $\left| f(x) - f(a) \right| >\varepsilon$. We can let $\delta = \frac{1}{n}$ and choose a sequence $x_n$. Now $x_n \to  a$, so by the assumption we have $f(x_n) \to f(a)$, but $\left| f(x_n) - f(a) \right| >\varepsilon$ for all $n$, a contradiction. Therefore $f$ must be continuous at $a$.
\end{proof}
\item For $f, g: [a,b] \to  \mathbb{R}$, show that $T(f \pm g) \le T(f) + T(g)$, and if $c \in \mathbb{R}$, then $T(cf) = |c| T(f)$.
\begin{proof}
	Fix some partition $x_0, \ldots, x_n$ of $[a,b]$. By triangle inequality we have for each $i = 1,\ldots, n$,
	\begin{align*}
		&\left| \left( f \pm g \right)(x_i) -  \left( f \pm g \right)(x_{i-1}) \right| \\
		&= 
\left| \left( f(x_i) - f(x_{i-1}) \right) \pm \left( g(x_i) - g(x_{i-1}) \right)  \right|\\
		&\le \left| f(x_i) - f(x_{i-1}) \right| + \left| g(x_i) - g(x_{i-1}) \right| 
	.\end{align*} 
	Passing to the sum, we have
	\[
		\sum_{i=1}^{n} \left( f \pm g \right) |_{x_{i-1}}^{x_i} 
		\le \sum_{i=1}^{n} f |_{x_{i-1}}^{x_i}
		+ \sum_{i=1}^{n} g |_{x_{i-1}}^{x_i}
	.\] 
	Taking supremum over the partition, we have
	\[
		T(f \pm g) \le  T(f) + T(g)
	.\] 
	To prove the second part of the proposition, fix some partition and observe that for each $i = 1,\ldots, n$,
	\[
		\left| (cf)(x_i) - (cf)(x_{i-1}) \right| 
		= |c| \left| f(x_i) - f(x_{i-1}) \right| 
	.\] 
	Now we can pass to sum and supremum and obtain the result.
\end{proof}
\item If $\varphi, \psi \in \text{BV}[a,b]$, does it follow that $\varphi \psi \in \text{BV}[a,b]$?
\begin{proof}
	Yes. We use that fact that BV functions are differences of positive increasing functions. Let $\varphi = \varphi_1 - \varphi_2$ and $\psi = \psi_1 - \psi_2$ where $\varphi_1, \varphi_2, \psi_1, \psi_2$ are positive increasing. Then
	\[
		\varphi\psi = (\varphi_1-\varphi_2)(\psi_1-\psi_2) = (\varphi_1\psi_1+\varphi_2\psi_2) - (\varphi_2 \psi_1 + \varphi_1 \psi_2)
	.\] 
	where the pairwise products of $\varphi_i \psi_i$ are also positive increasing.
\end{proof}
\item Show that $E \subset \mathbb{R}$ is of measure $0$ iff there exists collection of intervals $I_1, I_2, \ldots$ of finite total length such that each $x \in E$ is contained in infinitely many $I_j$.
\begin{proof}
	For each $n \in \mathbb{N}$, let $\{I^{(n)}_i\}_{i \in \mathbb{N}}$ be a collection of intervals that cover $E$ and with total length $\frac{\varepsilon}{2^n}$. Then for each point in $E$ and each $n$ there exists some $I^{(n)}_i$ containing $x$. Furthermore, the total length of all intervals is $\sum_{n=1}^{\infty} \frac{\varepsilon}{2^n} = \varepsilon$.
\end{proof}

	
\item Let us say that a number $c \in \mathbb{R}$ is unlucky if it can be written as a decimal without the digit $7$. Prove that the unlucky numbers form a set of measure $0$.
\begin{proof}
	We mimic the argument that Cantor set has measure zero.

	Let $F_0$ be the interval $[0,1]$. Let $F_1 = [0,0.7) \cup [0.8,1]$. Define $F_2$ as the set of all numbers whose second digit is not $7$. Proceed in this way, define
	\[
		F_n = \{x \in [0,1]: n\text{-th decimal digit of }x\neq 7\}  
	.\] 
	Let $c(X)$ denote the content of set $X$. We see that $c(F_n) = \frac{9}{10}$. Also, since each digits are independent of each other, we see that
	\[
		c\left(\cap_{n = 0}^{N}F_n\right) = \left( \frac{9}{10} \right) ^N
	.\] 
	Therefore
	\[
		c\left( \bigcap_{n=0} ^{\infty }F_n \right) =0
	.\] 
	Then note that, by definition of $F_n$, the unlucky numbers are exactly $\bigcap_{n=0} ^\infty F_n$.
\end{proof}



\end{itemize}
